\subsection{Implementation Details}

\subsubsection{Datasets}
\label{sec:datasets}
We conduct experiments on four real-world datasets spanning the short-video and movie domains; their statistics after preprocessing are summarized in Table~\ref{tab:dataset_stats}.

\begin{table*}[t]
\centering
\small
\caption{Statistics of the datasets after preprocessing. Categories are video categories, movie genres, or content tags depending on the dataset.}
\label{tab:dataset_stats}
\begin{tabular}{lccccc}
\toprule
Dataset & Users & Items & Categories & Interactions & Avg.\ interactions/user \\
\midrule
KuaiRec     & 7,176 & 10,728 & 31 & $\approx$9.57M & 1,334 \\
ML1M        & 6,034 & 3,125  & 18 & 574,376        & 95.2 \\
MicroLens   & 5,670 & 26,922 & 57 & 86,664         & 15.3 \\
KuaiRand-1K & 995   & 20,000 & 44 & 547,659        & 550.4 \\
\bottomrule
\end{tabular}
\end{table*}

\begin{itemize}
    \item \textbf{KuaiRec}~\cite{gao2022kuairec} is a large-scale real-world video recommendation dataset. The dataset contains 10,728 unique items across 31 semantic categories. Items and users are retained without explicit interaction thresholds (any positive exposure counts), but we apply a de-duplication pass that collapses repeated user--video pairs into a single occurrence, keeping the highest-watch-time instance. Following standard sequential recommendation protocols, we construct four dataset variants based on different item selection strategies for the train/validation/test split: \emph{highest\_individual}, \emph{highest\_average}, \emph{first\_individual}, and \emph{first\_average}. For each variant, we adopt leave-one-out: the last item in each user's chronological interaction sequence becomes the test item, the second-to-last item becomes the validation item, and the remaining items form the training set. All sequences are truncated to a maximum length of 50; shorter sequences are kept intact and padded with $0$ at the front.
    \item \textbf{MovieLens-1M (ML1M)}~\cite{harper2015} is a classic movie rating dataset. We keep ratings of \emph{at least 4 stars} as positive feedback and discard the rest; we then apply iterative 5-core filtering on users and items (a user is kept only if it has at least 5 positive interactions after item filtering, and vice versa) to remove long tails. Movies are described by 18 multi-label genres that serve as categories. Split protocol: chronological leave-one-out identical to KuaiRec — test = last item, validation = second-to-last, train = remainder, max sequence length 50.
    \item \textbf{MicroLens}~\cite{ni2023microlens} is a content-driven micro-video recommendation dataset with rich per-video category metadata. We retain only interactions whose \emph{watch-time-to-duration ratio is at least 0.1} (excluding near-zero-view events), then deduplicate user--video pairs by keeping the single highest-ratio engagement per pair. No further user/item core filtering is applied. Videos are organized into 57 content categories. Split and truncation protocol matches the other three datasets: chronological leave-one-out, max length 50.
    \item \textbf{KuaiRand-1K}~\cite{gao2022kuairand} is an unbiased short-video recommendation log collected under randomized exposure. We keep only \emph{fully watched} (\emph{long-view}) interactions, deduplicate user--video pairs, retain the 20{,}000 most popular videos by interaction count, and then apply iterative 5-core filtering on users and items. Videos carry 44 content tags. Split and truncation protocol: chronological leave-one-out, max length 50.
\end{itemize}

For all four datasets, items are re-indexed to $1,\ldots,|\mathcal{I}|$ with $0$ reserved for padding, and interactions are grouped by user and sorted chronologically before splitting. No dataset is merged or subsampled after the above preprocessing.

\subsubsection{Evaluation Protocol}
We adopt a \emph{two-tier evaluation} strategy to reconcile training efficiency with measurement fidelity.

\paragraph{Validation (early stopping).}
For every dataset, the held-out validation item is scored against a fixed pool of \emph{99 sampled negatives} drawn uniformly at random from the item catalog with \emph{seed 4444}. Every model on a given dataset uses the \emph{exact same} negative file (filename pattern
\verb|Dataset-random-sample_size=99-seed=4444.txt|), so validation numbers are comparable across models and checkpoint selection is deterministic. Validation scores are computed with a single held-out positive plus 99 negatives rather than the full catalog, purely for speed during training; this follows the standard sampled-protocol convention of early sequential recommenders~\cite{hidasi2016session,kang2018self}.

\paragraph{Test (final reporting).}
Final numbers reported in the paper are always \emph{full-catalog rankings}: each model produces a length-$k$ ranked list over \emph{all} $|\mathcal{I}|$ items and the held-out test item is checked for presence/rank against that full list, avoiding the well-known inconsistencies of sampled metrics~\cite{krichene2020sampled}. All models, including GRU4Rec, SASRec, BERT4Rec, DuoRec, TRIER, and PACER, are scored with the same shared full-ranking evaluator (function \texttt{evaluate\_function\_with\_full} in \texttt{script.py}) that takes only the positive item ID and the model's top-$k$ list; it never touches the negative file. This is consistent across all four datasets (KuaiRec variants, ML1M, MicroLens, KuaiRand-1K).

\paragraph{Metric definitions.}
We report accuracy metrics Recall@$k$ (Hit Ratio), NDCG@$k$, and MRR@$k$; diversity metrics Intra-List Distance (ILD@$k$)~\cite{ziegler2005} and Category Coverage (CC@$k$)~\cite{adomavicius2012}; and repetition-sensitive metrics Consecutive Similarity (CS@$k$) and MaxRun@$k$, for $k\in\{5,10,20\}$.

\noindent\textbf{ILD@$k$} is the average pairwise Euclidean distance between L2-normalized item category vectors:
\begin{equation}
\text{ILD@}k(R) = \frac{1}{k(k-1)}\sum_{s=1}^{k}\sum_{\substack{t=1\\t\neq s}}^{k}
                  \bigl\|\mathbf{c}_{j_s} - \mathbf{c}_{j_t}\bigr\|_2.
\end{equation}

\noindent\textbf{CC@$k$} is the fraction of distinct categories covered by list $R$:
\begin{equation}
\text{CC@}k(R) = \frac{\bigl|\bigcup_{s=1}^{k} \mathcal{T}(j_s)\bigr|}{C},
\end{equation}
where $C$ is the total number of categories in the dataset.

\noindent\textbf{CS@$k$} is the average cosine similarity between items at consecutive positions (lower CS means locally less repetitive lists):
\begin{equation}
\text{CS@}k(R) = \frac{1}{k-1}\sum_{s=2}^{k} \cos(\mathbf{v}_{j_s}, \mathbf{v}_{j_{s-1}}).
\end{equation}

\noindent\textbf{MaxRun@$k$} measures the longest contiguous segment of items that are mutually similar (same category \emph{or} high content-vector cosine), capturing the worst-case repetition streak where CS only reports the average. Let
\begin{equation}
a_s = \begin{cases}
  1 & \text{if } \mathcal{T}(j_s)\cap\mathcal{T}(j_{s-1})\neq\emptyset \;\text{or}\;
      \cos(\mathbf{v}_{j_s},\mathbf{v}_{j_{s-1}}) > \tau_{\text{run}},\\
  0 & \text{otherwise},
\end{cases}
\end{equation}
with $\tau_{\text{run}}=0.7$ throughout; then
\begin{equation}
\text{MaxRun@}k(R) = 1 + \max_{\text{subarray } i,\dots,j} \sum_{s=i+1}^{j} a_s,
\quad 1 \leq i < j \leq k,
\end{equation}
i.e., one plus the largest sum of $a_s$ over any contiguous window (a run of $L$ similar items contributes $L-1$ ones). A list with no adjacent similarities has MaxRun$=1$; a list where every adjacent pair is similar has MaxRun$=k$.

Crucially, ILD, CS, and MaxRun are measured for \emph{all} models in the same frozen, pretrained content-embedding space (the dataset-specific category/item2vec vectors released with the data), so cross-model diversity numbers are directly comparable.

\subsubsection{Model Configuration}

\paragraph{Training setup and common hyper-parameters.}
All models are implemented in PyTorch and share the same PyTorch \texttt{manual\_seed(42)}/\texttt{cuda.manual\_seed\_all(42)} seed for weight initialization and data ordering within a dataset. The maximum sequence length is 50 for every model, every dataset, no exceptions. We optimize with Adam~\cite{kingma2015}, batch size 256, learning rate $0.001$, for at most 1{,}000 epochs (early-stopping patience 100, min delta $10^{-4}$). Backbone-specific hyper-parameters (hidden size, number of layers, attention heads, dropout) are set to the values recommended in each baseline's original paper; we do not re-tune them on our data splits.

\paragraph{PACER-specific configuration.}
In addition to the shared backbone settings above, our PACER method has the following components and dedicated hyper-parameters. The initialisation of all learnable parameters follows a zero-mean Gaussian of standard deviation $0.02$. The embedding dimension is 64 and all linear mapping functions share the same hidden size. Both the RT and PT transformers consist of 2 layers with 2 attention heads, with dropout $0.5$ applied to the embedding matrix and the transformer module. The diversity trade-off $\lambda$ in Eq.~\eqref{eq:gen_score} is set to $0.1$ following~\cite{shi2024trier}, the inference-time position-aware penalty weight $\lambda_c$ to $0.01$, the soft order-temperature $\tau_o$ to $0.1$ (default), and the soft adjacent-order training weight $\gamma_o$ to $0.01$; the cosine margin in $\mathcal{L}_{order}$ is $m=0.5$. The number of retrospectively augmented items ($K$) is set to 5 and the number of augmented sequences, i.e., the beam width ($M$), to 3. The RT model is trained first to convergence and its parameters are frozen while the PT model is trained on top.

\paragraph{Checkpoint selection (no test leakage).}
For every model on every dataset, early stopping is driven by \emph{validation Recall@5} computed against the 99 fixed negative samples described above. The epoch that maximises \emph{validation} Recall@5 is saved as the best checkpoint and used for all test-time results. We do \emph{not} select checkpoints based on any test metric (including test HR@20, test NDCG@20, or any diversity metric), and we do not run a separate validation pass over the test set. This holds for every baseline as well as PACER, so the reported test numbers are unbiased estimates of generalisation.

\paragraph{Cross-model consistency.}
Every model (GRU4Rec, SASRec, BERT4Rec, DuoRec, TRIER, PACER) on a given dataset is trained with the exact same train/validation/test split, the same chronological ordering of interactions, the same random seed, and — at test time — the same full-ranking evaluator on the same candidate set. The only allowed difference across models is the architecture-specific hyper-parameter settings documented above. The diversity evaluator also reads from the same frozen content-embedding file across all models on a dataset, so ILD, CS, and MaxRun numbers are directly comparable and cannot be inflated by a model learning a different representation space.

\subsubsection{Baselines}
We compare with representative sequential recommendation methods, grouped into (i) vanilla sequential recommenders, (ii) self-supervised sequential recommenders, and (iii) diversified sequential recommenders.

\begin{itemize}
    \item \textbf{GRU4Rec}~\cite{hidasi2016session} is the pioneering recurrent approach to session-based recommendation. It embeds the item sequence into a stack of gated recurrent units (GRUs) that compresses the interaction history into a hidden state, which is then projected onto the item vocabulary to score candidates. We use one recurrent layer with embedding and hidden dimensions of 64 and dropout 0.2, trained with the standard next-item cross-entropy objective.
    \item \textbf{SASRec}~\cite{kang2018self} adapts the causal Transformer~\cite{vaswani2017} to sequential recommendation. It sums learned item and positional embeddings and applies a unidirectional self-attention mask so that each position can only attend to earlier items, predicting the next item at every sequence position in parallel. Following the original paper, we set the hidden size to 50 with 2 transformer blocks, 1 attention head, and dropout 0.5.
    \item \textbf{BERT4Rec}~\cite{sun2019bert4rec} replaces the causal mask with bidirectional self-attention and trains with a Cloze-style masked-item objective: a random subset of items (20\%) is replaced with a mask token and the model reconstructs them jointly using context from both sides, which strengthens contextual sequence modeling. We use a hidden size of 64, 2 transformer blocks, 2 attention heads, and dropout 0.2; masked training is converted to full-catalog next-item scoring at test time.
    \item \textbf{DuoRec}~\cite{wang2022duorec} is a self-supervised sequential recommender built on the same causal transformer backbone as SASRec. It regularizes representation learning with a contrastive loss between two views of a sequence: model-level augmentation obtained through stochastic dropout and \emph{supervised} positive pairs formed from sequences whose next target item is identical, which mitigates representation collapse. It is also the exact backbone architecture of our PT module ($d=64$, 2 layers, 2 heads), making it a clean reference for the effect of retrospective augmentation and diversity promotion.
    \item \textbf{TRIER}~\cite{shi2024trier} is the diversified sequential recommender most closely related to PACER and the foundation this work builds on. Its retrospective transformer (RT) processes the reversed interaction sequence and is trained with next-token prediction to hallucinate plausible items that could have preceded the observed history; beam search at inference yields $M$ augmented sequences, each interpreted as one of the user's latent intents. The prospective transformer (PT) then scores items by combining relevance with a diversity-promoting term that reserves probability mass for intents not yet covered by the current list, so diversity is optimized during list generation rather than by post-hoc re-ranking. We run the official implementation under our unified data splits and full-catalog ranking protocol, using plain item-ID embeddings and disabling content fusion, dense next-item supervision, and the position-aware consecutive penalty; any improvement of PACER over this configuration therefore isolates the contribution of our extensions.
\end{itemize}

\subsubsection{Ablation Variants}
To disentangle the contributions of PACER's components, we additionally report: (i)~\textbf{PT-no-ord}, which keeps the intent-coverage loss ($\lambda=0.1$) and the inference-time position-aware penalty ($\lambda_c=0.01$) but disables the soft adjacent-order training loss ($\gamma_o=0$), isolating the contribution of differentiable local-order supervision; (ii)~\textbf{PACER-full}, with all three diversity components enabled at $\lambda=0.1$, $\gamma_o=0.01$, $\lambda_c=0.01$; and (iii)~\textbf{RT-only}, the retrospective model evaluated alone for next-item prediction, assessing the standalone value of reverse trajectory modeling.
